\textbf{Proof.} Put \(r_b=b/(b+1)\), \(s_b^{\mathrm{sw}}=1/(b+1)\), \(q_b=r_b^{-b}=(1+1/b)^b\), and \(\lambda_b=r_b-s_b^{\mathrm{sw}}=(b-1)/(b+1)\). The assignment process is the stationary symmetric two-state Markov chain with transition matrix
\[
\begin{pmatrix}r_b&s_b^{\mathrm{sw}}\\s_b^{\mathrm{sw}}&r_b\end{pmatrix}.
\]
It is outcome-independent and therefore is an element of \(\mathcal D_T\). Stationarity and the Markov property give
\[
\upsilon_{t,b,a}=\frac12r_b^{k_t}.
\]
Since \(k_t\le b\),
\[
\upsilon_{t,b,a}\ge\frac12r_b^b=\frac{1}{2q_b}.
\]
The strict inequality \(q_b<e\) follows from \(b\log(1+1/b)<1\), so \(\upsilon_{t,b,a}>1/(2e)>0\). Conversely, for every \(b\), take \(T\ge b+1\) and \(t=b+1\). Then \(k_t=b\) and \(\upsilon_{t,b,a}=1/(2q_b)\). Since \(q_b\to e\), the uniform infimum is \(1/(2e)\), and it is not attained at any finite \(b\).

Fix \(Y\in\mathcal P_{T,\beta}(C)\), put \(y_{t,a}=Y_t(\mathbf a_{1:t})\), and define
\[
Z_t=\frac{H_{t,b,1}y_{t,1}}{\upsilon_{t,b,1}}-\frac{H_{t,b,0}y_{t,0}}{\upsilon_{t,b,0}},
\qquad
\widetilde\theta_b=\frac1T\sum_{t=1}^T Z_t.
\]
The two exposure events at time \(t\) are disjoint. Hence the inclusion-probability identity gives
\[
\mathbb E Z_t=y_{t,1}-y_{t,0},
\qquad
\mathbb E\widetilde\theta_b=\operatorname{GATE}_T(Y).
\]
Moreover, Assumption~\(\mathrm{ass:bounded-array}\) gives the sharper time-specific second-moment bound
\[
\mathbb E Z_t^2
=\frac{y_{t,1}^2}{\upsilon_{t,b,1}}+\frac{y_{t,0}^2}{\upsilon_{t,b,0}}
=2r_b^{-k_t}(y_{t,1}^2+y_{t,0}^2)
\le4B^2r_b^{-k_t}.
\]
If \(1\le u-t\le b\), Cauchy--Schwarz, \(\operatorname{Var}(Z_v)\le\mathbb E Z_v^2\), and the last display imply
\[
\left|\operatorname{Cov}(Z_t,Z_u)\right|
\le\sqrt{\operatorname{Var}(Z_t)\operatorname{Var}(Z_u)}
\le4B^2r_b^{-(k_t+k_u)/2}.
\]
For a longer separation \(u-t>b\), put \(d=u-b-t\ge1\). Conditional on \(W_{u-b}\), the following \(b\) coordinates remain equal to it with probability \(r_b^b\), and \(\upsilon_{u,b,a}=r_b^b/2\). Thus
\[
\mathbb E(Z_u\mid W_{u-b}=1)=2y_{u,1},
\qquad
\mathbb E(Z_u\mid W_{u-b}=0)=-2y_{u,0}.
\]
The transition identity for the symmetric chain is
\[
\Pr(W_{t+d}=j\mid W_t=i)
=\frac12\{1+(2i-1)(2j-1)\lambda_b^d\}.
\]
Consequently,
\[
\left|\mathbb E(Z_u\mid W_{1:t})-\mathbb EZ_u\right|
=\lambda_b^d|y_{u,1}+y_{u,0}|
\le2B\lambda_b^d.
\]
Since \(Z_t\) is measurable with respect to \(W_{1:t}\) and \(\mathbb E|Z_t|=|y_{t,1}|+|y_{t,0}|\le2B\),
\[
|\operatorname{Cov}(Z_t,Z_u)|\le4B^2\lambda_b^d.
\]
There are \((T-b-d)_+\) long-separation pairs with a given \(d\). Summing the diagonal bounds and twice the covariance bounds proves
\[
\operatorname{Var}\left(\sum_{t=1}^T Z_t\right)\le B^2V_{T,b}^{\sharp}.
\]
Unbiasedness and Jensen's inequality now yield
\[
\mathbb E\left|\widetilde\theta_b-\operatorname{GATE}_T(Y)\right|
\le\frac BT\sqrt{V_{T,b}^{\sharp}}.
\]
This is a finite variance certificate with exact pair counts, not an assertion that the upper bound equals the variance.

We compare it with the former coarse envelope
\[
V_{T,b}^{\mathrm{coarse}}
=4q_bT+8q_b\left(bT-\frac{b(b+1)}2\right)
+8\sum_{d=1}^{(T-b-1)_+}(T-b-d)\lambda_b^d.
\]
For every \(t\), \(r_b^{-k_t}\le q_b\), and for every short pair, \(r_b^{-(k_t+k_{t+h})/2}\le q_b\). Also
\[
\sum_{h=1}^b\sum_{t=1}^{T-h}1
=\sum_{h=1}^b(T-h)
=bT-\frac{b(b+1)}2.
\]
The diagonal inequality is strict at \(t=1\), because \(r_b^{-k_1}=1<q_b\). The long-lag terms coincide, so termwise comparison gives
\[
V_{T,b}^{\sharp}<V_{T,b}^{\mathrm{coarse}}.
\]
For completeness, we verify the uniform coarse bound. If \(b=1\), then \(q_b=2\), \(\lambda_b=0\), and \(V_{T,1}^{\mathrm{coarse}}=24T-16\le(8e+4)T\). If \(b=2\), then \(q_b=9/4\), \(\lambda_b=1/3\), and
\[
V_{T,2}^{\mathrm{coarse}}
\le45T-54+4(T-2)\le49T<(16e+8)T.
\]
For \(b\ge3\), the geometric tail is at most \((T-b)(b-1)/2\), whence
\[
V_{T,b}^{\mathrm{coarse}}
\le4q_bT+8q_b\left(bT-\frac{b(b+1)}2\right)+4(T-b)(b-1).
\]
The difference between \((8e+4)Tb\) and the right side is
\[
T\{8b(e-q_b)-4(q_b-1)\}+4b\{q_b(b+1)+b-1\}.
\]
Expanding \(q_b\) as
\[
q_b=\sum_{k=0}^b\frac1{k!}\prod_{j=0}^{k-1}\left(1-\frac jb\right)
\]
and retaining the deficits of the \(k=2\) and \(k=3\) terms gives \(e-q_b\ge b^{-1}-(3b^2)^{-1}\). Also \(q_b\le e<11/4\). Therefore
\[
8b(e-q_b)-4(q_b-1)
\ge8-\frac8{3b}-4(e-1)>0,
\]
and the remaining brace is positive. Hence
\[
V_{T,b}^{\sharp}<V_{T,b}^{\mathrm{coarse}}\le(8e+4)Tb.
\]
For fixed \(b\), only indices with \(k_t<b\) contribute to the diagonal difference, and only short pairs having at least one such index contribute to the short-pair difference. There are at most \(b\) such diagonal indices and at most \(b^2\) such pairs, while each summand is at most \(q_b<e\). Thus
\[
0<V_{T,b}^{\mathrm{coarse}}-V_{T,b}^{\sharp}=O(b^2)
\qquad(T\to\infty).
\]
Division of the coarse formula by \(Tb\) gives
\[
\lim_{T\to\infty}\frac{V_{T,b}^{\mathrm{coarse}}}{Tb}
=\frac{4q_b}{b}+8q_b+4\frac{b-1}{b}.
\]
The boundary difference vanishes after division by \(Tb\), so the same fixed-\(b\) limit holds for \(V_{T,b}^{\sharp}\). Since \(q_b\to e\),
\[
\lim_{b\to\infty}\lim_{T\to\infty}\frac{V_{T,b}^{\sharp}}{Tb}=8e+4.
\]
Thus \(8e+4\) remains the optimal coefficient for a uniform certificate \(V_{T,b}^{\sharp}\le KTb\).

It remains to compare observed outcomes with homogeneous-path outcomes. If \(t\le b+1\), then \(k_t=t-1\), so on \(H_{t,b,a}=1\) the complete history through \(t\) equals \(\mathbf a_{1:t}\); the replacement is exactly zero. If \(t\ge b+2\), Assumptions~\(\mathrm{ass:bounded-array}\) and \(\mathrm{ass:historical-polynomial-envelope}\) imply, on \(H_{t,b,a}=1\),
\[
|Y_t(W_{1:t})-y_{t,a}|
\le\min\left\{2B,C\sum_{\ell=b+1}^{t-1}(1+\ell)^{-(1+\beta)}\right\}.
\]
For each arm and time, \(\mathbb E[H_{t,b,a}/\upsilon_{t,b,a}]=1\). The triangle inequality therefore gives
\[
\mathbb E|\widehat\vartheta_b-\widetilde\theta_b|
\le\frac2T\sum_{t=b+2}^T
\min\left\{2B,C\sum_{\ell=b+1}^{t-1}(1+\ell)^{-(1+\beta)}\right\}
=E_{T,\beta}^{\sharp}(b).
\]
Dropping the boundedness cap and interchanging finite sums yields
\[
E_{T,\beta}^{\sharp}(b)
\le\frac{2C}{T}\sum_{\ell=b+1}^{T-1}(T-\ell)(1+\ell)^{-(1+\beta)}.
\]
Since \(x\mapsto x^{-(1+\beta)}\) is decreasing and \(\beta>0\), this is at most
\[
\frac{2C}{\beta}\frac{(T-b-1)_+}{T}(b+1)^{-\beta}
\le\frac{2C}{\beta}\frac{(T-b-1)_+}{T}b^{-\beta}.
\]
All sums are empty when their lower index exceeds the upper index. Thus replacement is zero through \(t=b+1\); this assertion includes \(b=1\), and total replacement is zero when \(T=1\) or \(b=T\).

Finally, the triangle inequality combines the sampling and replacement bounds to give
\[
\sup_{Y\in\mathcal P_{T,\beta}(C)}
\mathbb E|\widehat\vartheta_b-\operatorname{GATE}_T(Y)|
\le U_{T,\beta}^{\sharp}(b).
\]
Using \(V_{T,b}^{\sharp}\le(8e+4)Tb\) and the replacement-sum corollary gives
\[
U_{T,\beta}^{\sharp}(b)
\le\sqrt{8e+4}\,B\sqrt{\frac bT}
+\frac{2C}{\beta}\frac{(T-b-1)_+}{T}(b+1)^{-\beta}.
\]
The carryover quantity is an expected replacement-error upper bound obtained from the exposure-weight identity; no equality with a bias or maximal bias is claimed. \(\square\)